\section{Diversity, consensus, and population dynamics}
\label{sec:collective}
\subsection{Four phenomena often called an echo}
Similar wording, correlated errors, convergence during discussion, and synchronized action are distinct observations. Artificial Hivemind measures similarity among independent responses to open-ended prompts. Its reevaluation reproduces substantial within-model similarity while showing that embedding similarity has a high baseline even for different ideas on the same topic \citep{artificial-hivemind,artificial-hivemind-reevaluation}. The evidence supports repetition and concentration in that sampling setup, without measuring the effect of agent-to-agent exchange.

Error correlation instead requires questions with externally known answers. A panel may contain many samples yet repeatedly select the same wrong answer \citep{consensus-not-verification}. Under an illustrative exchangeable-error model with common variance and pairwise correlation $r$, the variance of an average yields the effective sample size
\begin{equation}
N_{\mathrm{eff}}=\frac{N}{1+(N-1)r}.
\end{equation}
This algebra explains why a large nominal panel need not supply many independent votes. It is not a fitted law of agent behavior: real error correlations depend on question, model, role, and interaction history. The Ringelmann-effect study directly estimates very small effective teams in dense debates among copies of small open models \citep{ringelmann-effect}; its reported 30-agent result should not be generalized to every modern team.

Convergence is a temporal property. Free conversations between model copies can approach repetitive styles or topics, with different endpoints across models and versions \citep{attractor-states-anthropic-card,attractor-states-jakkli-post,attractor-states-ko-geiping-paper,mythos-preview-self-interactions}. Long closed dialogues can expand their vocabulary while semantic measures plateau \citep{semantic-collapse}. These experiments deliberately remove external tasks and evidence. They isolate a self-reinforcing conversational tendency, while leaving open whether task feedback, tools, or changing environments prevent it.

Synchronization can occur without messages. In the road-choice experiment, identical policies responding to a common forecast can concentrate on one congested route; the effect varies across models and reasoning settings \citep{tokyo-road-game}. The shared environment supplies the coupling. Adding a communication channel is therefore neither necessary for a collective failure nor a sufficient explanation of one.

\subsection{Heterogeneity adds opportunities and selection problems}
Different models can contribute complementary correct answers, but a group must identify which participant is right. In the diversity-scaling experiments, small heterogeneous teams rival much larger homogeneous teams on average, with task-specific exceptions and a limited model range \citep{diversity-scaling}. Mixed three-model deliberation improves over the aggregate of the same models' independent answers in several settings, yet the strongest individual model remains better in most of the tested tasks \citep{deliberative-crowds}. These two baselines expose different benefits: interaction can improve a given pool without making that pool preferable to its best member.

A large heterogeneous pool can enlarge the oracle opportunity while worsening the implemented answer. The 23-model study reports that naive pooling and selection schemes can underperform the best individual even when at least one model often knows the answer \citep{mo-models-mo-problems}. Strong and weak models are not interchangeable samples. Majority voting can dilute expertise, and an imperfect model judge may favor familiar style over correctness. Heterogeneity should therefore be measured as complementarity conditional on capability and task, not simply as the number of vendors.

Interaction can consume the diversity that composition creates. In the interaction-tax experiments, agents exposed to complete peer solutions rapidly become more similar, and heterogeneous systems often lose relative to independent generation \citep{interaction-tax}. A scheme that keeps proposals separate until aggregation behaves differently. By contrast, SAT's learned strategy can use differences among models productively on reasoning tasks, but its homogeneous control executes strategies learned for the heterogeneous team \citep{sat-teams}. That asymmetry prevents a clean estimate of the intrinsic value of model diversity alone.

The weak-to-strong research campaign provides a practical example of diversity in search trajectories. A team given separate initial directions explores more successfully than a team given a shared prompt, whose active members converge onto one family of methods \citep{automated-w2s-researcher}. A later experiment does not reproduce the benefit of distributing starting ideas. The intervention is promising because it targets search coverage; its effectiveness is conditional on the task and how quickly subsequent communication erases the initial separation.

\subsection{Agreement can be a dynamical attractor}
Opinion studies distinguish a model's intrinsic preference from its response to social evidence. In a finite population, strong conformity and dense exposure can yield sharp changes in consensus as sampling noise varies \citep{okawa-biased-consensus}. A schematic control ratio in that work is
\begin{equation}
\kappa\,d/T,
\end{equation}
where $\kappa$ denotes conformity, $d$ the effective number of observed neighbors, and $T$ sampling temperature. We rename the conformity coefficient here to avoid confusing it with Ord's parallelization exponent. The ratio expresses a mechanism under the study's assumptions; it is not a universal threshold measurable from API temperature alone.

Exposure to others' choices produces herding. Following the design of the Music Lab experiment, Hashemi et al. (Cornell University and Stanford University) let large populations of LLM agents rate 48 newly generated songs, shown as titles and lyrics, with or without cumulative recommendation counts \citep{herding-runaway-cascades}. Exposure led to runaway cascades that diverged from the songs' intrinsic appeal and rarely corrected themselves. The authors set these outcomes beside the original human Music Lab, where people heard audio and saw download counts, as a qualitative benchmark rather than a controlled human--LLM comparison; most models herded more than the human groups, with model-dependent exceptions.

Private belief and public statement can also come apart, in the direction opposite to classic conformity, where people agree in public without changing their minds. Believed but Not Said (University of Illinois Chicago) asks each benign agent for its answer in a separate call that sees the discussion but whose reply is not shown to the group \citep{believed-but-not-said}. Averaged over six models and three datasets, the share of benign agents whose private answer at the final step is the false one rises to 44.8\% with one adversarial agent and 60.5\% with three. With three adversaries present, benign agents privately hold the false answer in 51.6\% of exposed turns but publicly endorse it in only 19.7\%. Much of the influence never appears in the transcript, so transcript-level agreement understates belief change.

Several observations limit a purely social explanation. A model can produce unanimous outputs because its own prior is strong, even if measured conformity is weak or negative. The pairwise Bayesian analysis separates attraction to a model's topic-specific position from influence exerted by a partner and finds the former often dominant \citep{brockers-opinion-dynamics}. Naming-game studies fit majority response curves and population transitions, but some large-population results use cached response policies rather than live LLM calls \citep{prev-demarzo-majority,prev-flint-group-size,indirect-tipping}. Those simulations test the extrapolation of a measured response rule; they do not show how a live model adapts across thousands of extended conversations. Passaro et al. extend this approach to networks, placing policies extracted from four LLMs on graphs of $10^3$ to $10^5$ agents playing the naming game \citep{network-topology-consensus}. Sparse networks slow consensus up to twelvefold. On the most heterogeneous topology, the median critical mass of committed agents needed to overturn an established convention is 3\% when they occupy hubs, 16\% when placed at random, and 35\% at low-degree nodes.

The Flag Game measures correctness as well as group state. Accuracy peaks around an intermediate group size in its principal configuration, while failure shifts from unanimous error toward division into competing camps as groups grow \citep{flag-game}. The peak is small relative to the reported dispersion, and group size is varied mainly on one model. Nevertheless, the distinction between correct consensus, incorrect consensus, and polarization is a useful reporting pattern. A single agreement score would conceal the change in failure type.

Group success and efficiency are also separate outcomes. In a collective-risk dilemma, six players contribute over ten rounds to a pool that protects their remaining resources from catastrophe \citep{cooperation-without-adaptation}. Four of the five tested models contribute about the fair share at every positive risk level, and a different set of four also contributes when the risk is zero, which helps nobody; stating explicitly that only the agent's own payoff matters reduces these zero-risk contributions for most models. In self-play, Gemini 3.5 Flash-Lite and Grok 4.20 both reach the group target, but their expected payoff per seat is 20.0 versus 5.0, because Grok keeps contributing after the target is met. Cooperative behavior and group success do not establish adaptation to incentives.

Piao et al.'s networked-population study, now published in Nature Communications, adds a larger-scale setting for polarization \citep{piao-polarization}. We treat its journal version and earlier preprint as one evidence lineage. Network structure, information presentation, and initial agent settings belong alongside model identity when interpreting such results.

RHEON supplies a live-inference comparison \citep{rheon-consensus}. Homogeneous frozen Qwen3-14B populations of eight to 1,000 agents repeatedly rewrite answers after reading their neighbors and their own previous answer. Semantic cohesion rises broadly over ten sweeps; factual improvement is less uniform. After excluding one false-premise question whose scoring was inverted, average factual repair exceeds average factual damage on the five retained questions. Even so, 85 of 1,800 runs worsen the hallucination fraction by more than 0.1. Geometry and population size partly change together; overlapping 64-agent settings permit narrower topology comparisons. The reported 4.77 million responses are inference events, not distinct agents.

The evaluation is part of this result. RHEON's Qwen2.5-72B judge receives the question and response without a supplied gold answer. Human validation revealed the inversion that prompted the exclusion. This makes the study useful both for separating semantic agreement from factual correctness and for showing how a judge can misidentify their relationship. Its limited question set and corrected scoring do not establish a general accuracy benefit from consensus.

A 2 October preprint clarifies the update mechanism. \citet{coherence-driven-contagion} elicit adoption from ADOPT/DECLINE first-token logits after prompting a model with the number of endorsing peers. Population dynamics then use these measured kernels in 30-agent simulations with Qwen2.5-3B-Instruct and Mistral-7B-Instruct-v0.3. In one matched-degree-six condition, a clustered ring reaches a final believing fraction of 0.67 versus 0.43 for the rewired/random network, across 100 replicates. This is a conditional topology effect on belief propagation, not a measure of task-solving quality. Its hysteresis experiments depend on a specified pressure ramp and separately measured retention rule. The small models, stylized domains, and kernel-driven updates limit transfer to free-form frontier-agent interaction.

\subsection{Intervention and measurement inside a collective}
Behavioral interventions can replace an agent's observation or message and track the downstream change, as in the Flag Game. Activation interventions can ask whether a particular internal representation mediates a response to social information. Social Circuits recovers message-induced changes by patching activations in small open models, with effects that vary strongly across tasks \citep{social-circuits}. The channel-assignment study likewise finds substantial susceptibility differences when the same content appears as user, assistant, or tool text \citep{not-just-rlhf-activations}. These results motivate mechanistic analysis without justifying a claim that frontier models possess a distinct ``agent-to-agent circuit.'' Behavioral probes offer a cheaper per-interaction measure: scores originally built for hallucination analysis estimate how much a node relies on external context versus its parametric knowledge, which tracks information flow in swarms of prompt-defined experts \citep{behavioral-probes-information-flow}.

Interventions triggered by behavioral signals can trade off prevention against useful correction. In homogeneous-panel experiments, freezing the debate on such a signal averts some transitions from a correct majority to a wrong answer while preventing more transitions in the beneficial direction \citep{homogeneous-panel-collapse}. Counting only averted failures would misstate the net outcome. HELIOS/NARCBench detects instructed collusion well in its tested settings, while transfer to less scripted behavior is much weaker \citep{narcbench-helios}. Capability to detect a constructed signal and readiness for general deployment are different claims.

\subsection{Simulation validity and observational populations}
The PIMMUR audit and robustness studies show that experiment description, persona format, memory, and knowledge of the hypothesis can materially alter simulated social outcomes \citep{pimmur-principles,prev-robustness-audits}. These are not generic disclaimers: they identify manipulable causes that can compete with the purported collective mechanism. Robustness should be tested through alternative prompt formulations, concealed hypotheses, and held-out environments before a numerical tipping point is treated as stable.

Field observations remove some artificial constraints while introducing different confounds. The German-wiki studies find rapid copying of recent pages and conventions, but names do not map one-to-one to runs and task success is only partially observed \citep{wiki-copying,lutje-mechanics-of-a-swarm}. Moltbook studies disagree partly because they measure different objects: vocabulary associated with institutions, reciprocal interaction, embedding movement, or inferred human control \citep{zhang-agents-in-the-wild,moltbook-socialization,moltbook-illusion,lai-looking-for-something-weird}. These observations can coexist. Institutional language need not imply durable authority; a stable daily embedding need not imply that every post is similar. The longest record now covers Moltbook from 27 January to 6 October 2026: 4,189,979 posts and 19,617,496 comments from 184,909 active accounts \citep{moltbook-collective-behavior}. Approval grows sublinearly with discussion size, with an exponent near 0.79. Coverage of new posts falls to about 85\% after June, and the discussion-tree analyses exclude posts with at least 100 comments, which are 0.49\% of posts but about 62\% of comments, so typical threads say little about the most viral ones.

The Superminds Test further separates participation from reasoning. Most posted questions receive no answer, so a low platform-level success rate is largely a failure to recruit work rather than a direct test of what millions of cooperating agents could solve \citep{superminds-test}. AI Village supplies persistent identities and known model assignments, but goals and interventions come from human operators and there is one evolving population \citep{ai-village-how-it-works}. EinsteinArena provides externally checked mathematical submissions and traces of reuse, while still lacking matched resource controls for collective versus individual search \citep{einsteinarena}. Each setting offers a useful observation when its unit of analysis is preserved.
